\chapter{什么时候应该使用 Agent}

\begin{takeaway}[学习目标]
完成本章后，你应该能区分聊天、工作流与 Agent；能用任务不确定性、环境反馈和风险判断是否需要 Agent；能为一个候选场景写出可测量的完成标准。
\end{takeaway}

\section{Agent 不是“更长的 Prompt”}

单次模型调用接收输入并返回输出。Agent 则处在一个循环中：它根据当前状态选择动作，通过工具改变或观察外部环境，再根据结果继续行动。Anthropic 将预定义代码路径称为 workflow，将由模型动态决定步骤和工具的系统称为 agentic system 中更自主的一类\citep{anthropic_effective_agents}。

\begin{definitionbox}{本书中的 Agent 定义}
Agent 是一个以模型为决策组件、能够在多步循环中观察环境、选择工具、更新状态，并以可检查的终止条件完成任务的软件系统。
\end{definitionbox}

这个定义刻意包含五个约束：多步、环境、工具、状态、终止。没有这些约束，“Agent”很容易退化成营销名称。

\section{复杂度阶梯}

面对一个需求，优先选择能完成任务的最低复杂度方案：

\begin{enumerate}
  \item \textbf{确定性代码}：规则明确、输入结构稳定，普通程序最好。
  \item \textbf{单次模型调用}：需要语言理解或生成，但不需要多步环境交互。
  \item \textbf{增强型调用}：模型加检索、结构化输出或一个受控工具。
  \item \textbf{工作流}：步骤可以预先定义，需要路由、并行或质量关卡。
  \item \textbf{Agent}：具体步骤无法预知，需要根据真实反馈动态决策。
  \item \textbf{多 Agent}：任务存在可证明的并行、隔离或专业化收益。
\end{enumerate}

每上升一级，你通常都要付出更多延迟、成本、调试难度和安全风险。复杂度不是能力勋章，而是一笔需要回报的投资。

\section{三个关键判断}

\subsection{步骤能否提前写死}

如果所有步骤都能在白板上稳定画出，优先工作流。例如“读取发票、提取字段、校验金额、写入数据库”可以有明确路径。相反，“调查某个市场为什么突然变化”可能需要根据每轮发现决定下一步搜索什么。

\subsection{环境是否提供真实反馈}

Agent 最擅长有反馈的环境：代码可以运行测试，网页会返回状态，数据库会返回查询结果，文件系统会告诉你文件是否存在。环境反馈把模型从“继续猜”拉回“观察事实”。

纯开放式写作缺少客观反馈，盲目循环常常只是让模型反复改写；此时 evaluator、人工编辑或明确 rubric 比自主循环更重要。

\subsection{错误是否可以被控制}

读取公开网页和删除生产数据库不是同一类工具。任务越自主，错误传播距离越长。选择 Agent 前必须知道：错误能否检测？能否撤销？是否需要人工确认？最坏损失是什么？

\begin{warningbox}
当一个动作不可逆、涉及资金、权限、隐私、医疗或法律后果时，不要把“模型认为可以执行”当作授权。把确认点写进系统，而不是写在团队约定里。
\end{warningbox}

\section{Agent 适用性评分表}

给每项打 0--2 分：0 表示不符合，2 表示高度符合。

\begin{table}[htbp]
\centering
\begin{tabularx}{\textwidth}{>{\bfseries}p{0.24\textwidth} X p{0.15\textwidth}}
\toprule
维度 & 判断问题 & 分数 \\
\midrule
步骤不确定性 & 是否必须根据中间结果决定下一步？ & 0--2 \\
工具需求 & 是否需要多个外部系统或工具？ & 0--2 \\
环境反馈 & 是否能从执行结果得到事实信号？ & 0--2 \\
任务价值 & 自动完成一次任务能否节省显著时间或成本？ & 0--2 \\
错误可控性 & 错误是否容易检测、隔离和恢复？ & 0--2 \\
评测可行性 & 是否能构造代表性任务集和完成标准？ & 0--2 \\
\bottomrule
\end{tabularx}
\end{table}

总分低于 6，通常不需要 Agent；6--9 可以先做受控工作流；10 分以上值得做 Agent 原型。评分不是科学定律，它的价值在于强迫团队讨论隐含假设。

\section{把“帮我做”改写成任务合同}

模糊需求：“帮我研究竞品。”

可执行任务合同：

\begin{codeblock}
目标：比较 5 个指定竞品最近 12 个月的核心功能与定价。
输入：竞品名单、目标用户、允许访问的公开来源。
输出：结构化表格、每个结论的来源、3 条产品机会。
完成条件：5 个竞品均覆盖；关键字段完整率 >= 90%；
          每个事实至少有一个原始来源。
停止条件：连续 3 次检索没有新增可靠信息；达到预算上限。
人工确认：发布报告或向外部联系人发消息之前。
\end{codeblock}

任务合同让系统知道何时停，也让评测者知道什么叫成功。

\begin{practice}[为一个场景做技术降级]
选择你想做的 Agent，分别写出四种实现：普通代码、单次模型调用、固定工作流、Agent。比较每种方案的成本、延迟、可控性和预期完成率。只有 Agent 版本有明确增量时才保留它。
\end{practice}

\section{常见误判}

\begin{itemize}
  \item \textbf{把聊天长度当自主性}：多轮聊天不等于系统能观察和行动。
  \item \textbf{把工具数量当能力}：十个描述含糊的工具不如两个契约清晰的工具。
  \item \textbf{把角色扮演当多 Agent}：不同名字的 prompt 不会自动产生专业化收益。
  \item \textbf{把 Demo 成功当生产可靠}：一次成功无法说明边界、成本和稳定性。
  \item \textbf{把模型升级当系统修复}：工具、状态和验收设计的问题不会因模型更大自动消失。
\end{itemize}

\begin{checkpoint}
\begin{itemize}
  \item \checkbox 我能用一句话区分 workflow 与 Agent。
  \item \checkbox 我为候选场景完成了六维评分。
  \item \checkbox 我写出了输入、输出、完成、停止和人工确认条件。
  \item \checkbox 我能解释为什么不用更简单的方案。
\end{itemize}
\end{checkpoint}

